% SPDX-FileCopyrightText: Copyright (c) 2016-2026 Objectionary.com
% SPDX-License-Identifier: MIT

\section{Semantics via \texorpdfstring{\(\varphi\)-Calculus}{phi-Calculus}}\label{sec:semantics}

The meaning of an EO program is given by translating it to a
  \(\varphi\)-calculus expression.
The \(\varphi\)-calculus is defined elsewhere~\citep{bugayenko2021eolang};
  the present section assumes its syntax and reduction rules
  (formation, application, dispatch, locator, data, function, dataization)
  and specifies the translation function
\begin{equation*}
  \llbracket \cdot \rrbracket\,:\,
    \text{well-formed EO expressions}\,\longrightarrow\,
    \varphi\text{-expressions}.
\end{equation*}
The translation is defined construct-by-construct over the outer kinds of
  \cref{app:outer-kinds}, using \cref{sec:lexical,sec:syntax} as the
  domain and the \(\varphi\)-calculus as the codomain.

\subsection{Translation Schema}\label{sec:sem:schema}

In equations that show how a rule expands, we write $\llbracket e \rrbracket$
  whenever $e$ is a sub-expression that is itself translated by the
  rules of this chapter.
In worked examples (\cref{app:examples} and the equations within this
  chapter), when the operand is a bare identifier or literal whose
  translation is the identity, we drop the brackets and write the
  identifier directly --- $\mathtt{foo}\,(\alpha_0 \mapsto \mathtt{a})$
  rather than $\llbracket \mathtt{foo} \rrbracket\,(\alpha_0 \mapsto
  \llbracket \mathtt{a} \rrbracket)$.

\Cref{tab:translation} summarises the translation per outer kind;
  the subsections that follow expand each row in turn.
A program in source is a top-level object whose translation
  \(\llbracket P \rrbracket\) is attached to \(\Phi\), the global root:
\begin{equation*}
  \Phi\,\mapsto\,\llbracket P \rrbracket.
\end{equation*}

\begin{table*}
\capt{Translation \(\llbracket \cdot \rrbracket\) per outer kind.}
\label{tab:translation}
\small
\begin{tabularx}{\textwidth}{l X}
\toprule
Outer kind & \(\varphi\)-translation \\
\midrule
\ff{head} (identifier or literal) &
  Locator dispatch (\cref{sec:sem:dispatch}) or data (\cref{sec:sem:atoms}). \\
\ff{hmethod}, \ff{vmethod} (no hargs) &
  Repeated dispatch: \(\llbracket h.\mathtt{m_1}.\mathtt{m_2}\ldots\rrbracket = \llbracket h \rrbracket.\mathtt{m_1}.\mathtt{m_2}\ldots\). \\
\ff{happlication}, \ff{vapplication}, \ff{vmethod-with-hargs} &
  Application of the head to its arguments: \(\llbracket h\,e_1\,e_2\,\ldots\rrbracket = \llbracket h \rrbracket\,(\alpha_0 \mapsto \llbracket e_1 \rrbracket, \alpha_1 \mapsto \llbracket e_2 \rrbracket, \ldots)\), with binding labels in place of \(\alpha_i\) when explicit. \\
\ff{bare-formation} &
  Formation \(\llbracket B \rrbracket\) collecting each child as a pair (\cref{sec:sem:formation}). \\
\ff{bare-reversed}, \ff{reversed-with-hargs} &
  Dispatch on the receiver: \(\llbracket \mathtt{m.}\,r\,e_1\,\ldots \rrbracket = \llbracket r \rrbracket.\mathtt{m}\,(\alpha_0 \mapsto \llbracket e_1 \rrbracket, \ldots)\). \\
\ff{compact-tuple} &
  Application with the trailing children wrapped in \(\Phi.\mathtt{tuple}\) (\cref{sec:sem:tuple}). \\
\ff{inline-phi-formation} &
  A fresh formation whose \(\varphi\) attribute is the LHS body (\cref{sec:sem:inline-phi}). \\
\ff{text-block} &
  A string literal whose \(\Delta\)-asset is the UTF-8 bytes of the block body. \\
\bottomrule
\end{tabularx}
\end{table*}

\subsection{Formations and Voids}\label{sec:sem:formation}

A formation \ff{[a b] > foo} translates to a binding
\begin{equation*}
  \mathtt{foo}\;\mapsto\;
    \llbracket\,
      a \mapsto \varnothing,\ b \mapsto \varnothing
    \,\rrbracket
\end{equation*}
in the enclosing scope.
Each name in the parameter list contributes a void pair.
The \ff{@} parameter is admitted and translates to the special name
  \(\varphi\); \ff{\textasciicircum} is forbidden as a parameter name
  (\cref{sec:syntax:formation}).

A formation with a body block contributes each child as an additional
  pair.
A child whose name suffix is \ff{> name} contributes the pair
  \(\mathtt{name} \mapsto \llbracket e \rrbracket\) where \(e\) is the
  child expression.
A child whose suffix is \ff{>>} contributes a pair under the
  auto-generated name (\cref{sec:syntax:names});
  the auto-name is treated as an ordinary attribute label in the
  translation.
A child whose suffix is \ff{+> name} contributes a \emph{test attribute}
  (see \cref{sec:wf:atoms}); it is represented in the codomain as an
  ordinary pair whose key is \(\mathtt{+name}\)
  --- i.e., the user-supplied \ff{name} with a literal \ff{+} prefix.
A child with no suffix (a master child --- a nested formation, atom, or
  inline-phi formation) contributes a pair whose key is the master child's
  own name.

\subsection{Application}\label{sec:sem:application}

Throughout the translation, the labels
  \(\alpha_0, \alpha_1, \ldots\)
  denote successive positional attribute slots of a formation, in the
  convention of the \(\varphi\)-calculus.
They name the slot a positional argument fills, not the argument value
  itself.

A horizontal application \ff{h e\(_1\) e\(_2\) \ldots} translates to
\begin{equation*}
  \llbracket h\,e_1\,e_2\,\ldots \rrbracket\;=\;
    \llbracket h \rrbracket\,
    (\,
      \alpha_0 \mapsto \llbracket e_1 \rrbracket,\
      \alpha_1 \mapsto \llbracket e_2 \rrbracket,\
      \ldots
    \,).
\end{equation*}
A vertical application of the same arity translates identically:
  positional indices are read from the source order of the deeper-indent
  children.

Inline bindings replace positional indices with explicit attribute
  labels.
The argument \ff{e:label} translates to the pair
  \(\mathtt{label} \mapsto \llbracket e \rrbracket\);
  the argument \ff{e:N} translates to the pair
  \(\alpha_N \mapsto \llbracket e \rrbracket\).
The well-formedness rules of \cref{sec:wf:bindings} guarantee that an
  argument group is either uniformly positional or uniformly labelled
  --- the translation never mixes \(\alpha_i\) labels with named pairs in
  one application.

An application that fills only some voids of its head is a \emph{partial
  application}.
Its result is a fresh formation whose voids are the head's voids
  \emph{minus} the filled ones, with the filled ones replaced by their
  argument expressions in the binding list.
For instance, if \ff{foo} is \(\llbracket\,a \mapsto \varnothing,\ b \mapsto \varnothing\,\rrbracket\),
  then
  \(\llbracket \mathtt{foo}\,e \rrbracket = \mathtt{foo}\,(a \mapsto \llbracket e \rrbracket)\)
  reduces to
  \(\llbracket\,a \mapsto \llbracket e \rrbracket,\ b \mapsto \varnothing\,\rrbracket\) ---
  an abstract formation with one void remaining.
The full reduction system is the \(\varphi\)-calculus's;
  \citet{bugayenko2021eolang} gives the rule.

\subsection{Dispatch and Locators}\label{sec:sem:dispatch}

A dispatch \(e.\mathtt{a}\) translates to a \(\varphi\)-dispatch
  \(\llbracket e \rrbracket.\mathtt{a}\).
Method chains expand left-to-right:
  \(\llbracket e.\mathtt{a}_1.\mathtt{a}_2\ldots\rrbracket
    = \llbracket e \rrbracket.\mathtt{a}_1.\mathtt{a}_2\ldots\).
Reversed dispatch \(\mathtt{m.}\,r\,e_1\,\ldots\) translates to the
  postfix form
  \(\llbracket r \rrbracket.\mathtt{m}\,(\alpha_0 \mapsto \llbracket e_1 \rrbracket,\ldots)\):
  reversed and postfix dispatch denote the same \(\varphi\)-term
  (\cref{sec:over:reversed}).

The four single-character locator tokens of \cref{sec:lexical:tokens}
  are translated to the corresponding \(\varphi\)-calculus names:

\begin{center}
\begin{tabularx}{0.7\textwidth}{l X l}
\toprule
EO token & Role & \(\varphi\)-name \\
\midrule
\ff{Q} (\ff{ROOT}) & The global root namespace            & \(\Phi\) \\
\ff{\$} (\ff{XI})  & The current scope (self)             & \(\xi\)  \\
\ff{@} (\ff{PHI})  & The decoratee of the enclosing object & \(\varphi\) \\
\ff{\textasciicircum} (\ff{RHO}) & The parent of the enclosing object & \(\rho\) \\
\bottomrule
\end{tabularx}
\end{center}

The locator \(\xi\) is the implicit prefix of any unqualified dispatch:
  \(\llbracket .\mathtt{a} \rrbracket = \xi.\mathtt{a}\).
The locator \(\rho\) refers to the formation one level outward;
  it has no source-level counterpart at indent 0 (a top-level formation
  has no parent).
The locator \(\Phi\) anchors an absolute path:
  \(\llbracket \mathtt{Q.io.stdout} \rrbracket = \Phi.\mathtt{io}.\mathtt{stdout}\).
The locator \(\varphi\) is the \emph{decoratee}: a dispatch on a formation
  that finds no local attribute falls through to the same dispatch on
  the formation attached to \(\varphi\).
This fall-through is the entire mechanism of behavioural reuse in EO ---
  there is no class hierarchy and no implementation inheritance.

\subsection{Assets, Data, and Atoms}\label{sec:sem:atoms}

A formation may, in addition to its ordinary attribute pairs, carry one
  or both of two reserved \emph{assets}:
\begin{itemize}
  \item a \emph{$\Delta$-asset}, holding a finite sequence of 8-bit
    bytes (raw data);
  \item a \emph{$\lambda$-asset}, holding a host-language function
    (a hook into a function implemented outside EO).
\end{itemize}
An asset is keyed by the reserved name $\Delta$ or $\lambda$ and does
  not occupy an ordinary attribute slot.
Assets are not introduced by any surface construct directly;
  they are produced by the translation of data literals and atoms,
  described below.

\paragraph{Data primitives.}
Data literals translate to formations whose $\Delta$-asset carries a
  fixed byte encoding of the literal value, decorated by the
  appropriate canonical object.
The encodings, by lexical class, are:

\begin{itemize}
  \item \ff{INT} --- 8 bytes, big-endian, two's-complement signed
    64-bit integer.
    Decorated by $\Phi.\mathtt{number}$.
  \item \ff{FLOAT} --- 8 bytes, big-endian, IEEE-754 binary64
    (double-precision) floating point.
    Decorated by $\Phi.\mathtt{number}$.
  \item \ff{HEX} --- the same 8-byte two's-complement integer
    encoding as \ff{INT}, after the hex digits are interpreted as a
    signed 64-bit integer.
    Decorated by $\Phi.\mathtt{number}$.
  \item \ff{STRING} --- the UTF-8 byte sequence of the string content
    after escape sequences are resolved (\cref{sec:lexical:strings}).
    Decorated by $\Phi.\mathtt{string}$.
  \item \ff{TEXT} --- the UTF-8 byte sequence of the text block's
    body, joined by \(\backslash\)n with the opening indent stripped
    and with escape sequences resolved.
    Decorated by $\Phi.\mathtt{string}$.
  \item \ff{BYTES} --- the byte sequence written literally, in order.
    Decorated by $\Phi.\mathtt{bytes}$.
\end{itemize}

For example, the integer literal \ff{42} translates to
  $\Phi.\mathtt{number}\,(\Phi.\mathtt{bytes}\,(\,
    \llbracket \Delta \mapsto \mathtt{40{-}45{-}00{-}00{-}00{-}00{-}00{-}00} \rrbracket
  \,))$,
  matching the encoding given by \citet{bugayenko2021eolang}.

\paragraph{Tuples.}
The canonical object $\Phi.\mathtt{tuple}$ is a single recursive cons-cell:
  a formation with three voids $\mathtt{tail}$, $\mathtt{head}$, and
  $\mathtt{length}$.
A non-empty tuple is built by applying $\Phi.\mathtt{tuple}$ with
  $\mathtt{tail}$ bound to a shorter tuple, $\mathtt{head}$ bound to the
  new element, and $\mathtt{length}$ bound to the resulting length.
The empty tuple is the canonical object $\Phi.\mathtt{tuple}.\mathtt{empty}$.
The $\mathtt{length}$ void is set explicitly, not computed at dispatch
  time, so that $\mathtt{length}$ has constant-time cost.

A surface form \ff{* e\(_1\) e\(_2\) \ldots e\(_n\)} unrolls into the
  right-associated chain
\begin{equation*}
  \Phi.\mathtt{tuple}\,(\,
    \mathtt{tail} \mapsto \cdots\Phi.\mathtt{tuple}\,(\,
      \mathtt{tail} \mapsto \Phi.\mathtt{tuple}.\mathtt{empty},\
      \mathtt{head} \mapsto \llbracket e_1 \rrbracket,\
      \mathtt{length} \mapsto 1
    \,)\cdots,\
    \mathtt{head} \mapsto \llbracket e_n \rrbracket,\
    \mathtt{length} \mapsto n
  \,).
\end{equation*}
Concretely, \ff{* x y z} translates to
\begin{equation*}
  \Phi.\mathtt{tuple}\,(\,
    \mathtt{tail} \mapsto \Phi.\mathtt{tuple}\,(\,
      \mathtt{tail} \mapsto \Phi.\mathtt{tuple}\,(\,
        \mathtt{tail} \mapsto \Phi.\mathtt{tuple}.\mathtt{empty},\
        \mathtt{head} \mapsto \mathtt{x},\
        \mathtt{length} \mapsto 1
      \,),\
      \mathtt{head} \mapsto \mathtt{y},\
      \mathtt{length} \mapsto 2
    \,),\
    \mathtt{head} \mapsto \mathtt{z},\
    \mathtt{length} \mapsto 3
  \,).
\end{equation*}
The standard library provides $\mathtt{at}$ (indexed access),
  $\mathtt{length}$ (the void retrieved as-is), and similar canonical
  attributes on $\Phi.\mathtt{tuple}$ that walk the chain;
  these are ordinary canonical attributes and are not part of the
  language proper.
We write $\Phi.\mathtt{tuple}(e_1, \ldots, e_n)$ as shorthand for the
  unrolled chain above when the internal structure is not the point of
  the discussion (e.g.\ in \cref{sec:sem:tuple}).

\paragraph{Atoms.}
An atom is a formation whose name suffix is \ff{> name /sig}.
The signature \ff{sig} declares the atom's return type;
  the operational behaviour is supplied by a host-language function
  attached to the formation's $\lambda$-asset:
\begin{equation*}
  \llbracket \mathtt{[x]\ >\ plus\ /number} \rrbracket\;=\;
    \mathtt{plus}\,\mapsto\,
    \llbracket\,
      x \mapsto \varnothing,\
      \lambda \mapsto \mathtt{Plus}
    \,\rrbracket
\end{equation*}
where $\mathtt{Plus}$ is a host-language function whose return matches
  the signature.
An atom's body may contain only test attributes (\cref{sec:wf:atoms}),
  which translate to ordinary $\mathtt{+name}$ pairs.

\paragraph{Dataization.}
Evaluating an EO program reduces it to a sequence of bytes via the
  \(\varphi\)-calculus dataization operator
\begin{equation*}
  \mathbb{D}\,:\,\varphi\text{-expression}\,\to\,\text{bytes},
\end{equation*}
which normalises its argument and, if the normal form carries a
  $\Delta$-asset, returns those bytes;
  otherwise it invokes the $\lambda$-asset of the normal form and
  dataizes its result recursively.
Dataizing the top-level program $\Phi$ is the evaluation of the program.

\subsection{Compact Tuple}\label{sec:sem:tuple}

The compact-tuple line \ff{head *N} (\cref{sec:syntax:compact})
  with vertical children \(c_1, \ldots, c_K\) translates to
\begin{equation*}
  \llbracket h \rrbracket\,(
    \alpha_0 \mapsto \llbracket c_1 \rrbracket,\ \ldots,\
    \alpha_{N-1} \mapsto \llbracket c_N \rrbracket,\
    \alpha_N \mapsto \Phi.\mathtt{tuple}(
      \llbracket c_{N+1} \rrbracket,\ \ldots,\ \llbracket c_K \rrbracket
    )
  ).
\end{equation*}
The first \(N\) children become direct positional arguments of the head;
  the remaining \(K - N\) children fold into a single
  \(\Phi.\mathtt{tuple}\) argument that takes the trailing positional
  slot.
The construct is pure syntactic sugar over the application form of
  \cref{sec:sem:application};
  it is included in the surface because tupling trailing arguments is
  pervasive in practice
  (e.g.\ argument lists of \ff{sprintf}, sequences of statements in
  \ff{seq}).

A bare \ff{*} appearing in head position is itself a
  \(\Phi.\mathtt{tuple}\) constructor: an application whose head is \ff{*}
  with horizontal arguments translates to
  \(\Phi.\mathtt{tuple}(\llbracket e_1 \rrbracket, \ldots, \llbracket e_n \rrbracket)\).

\subsection{Const-marker}\label{sec:sem:const}

A name suffix that ends in \ff{!} (\ff{> name!} or \ff{>>!}) marks the
  attribute as a \emph{constant}: once bound, the binding may not be
  rebound by application.
Const-marking does not change the \(\varphi\)-translation of the formation
  itself --- the pair $\mathtt{name} \mapsto e$ is emitted unchanged.
It is a hint to downstream tooling that the binding is final and may be
  cached or rewritten under the assumption that it will not change.

\subsection{Inline-phi Formation}\label{sec:sem:inline-phi}

A line whose suffix is \ff{> [params] > name} introduces an
  \emph{inline-phi formation}.
The left-hand expression of the suffix becomes the body of a fresh
  formation, bound to its \(\varphi\) attribute, and the right-hand name
  binds the formation in the enclosing scope:
\begin{equation*}
  \llbracket \mathtt{e\ >\ [p_1\ p_2\ \ldots]\ >\ name} \rrbracket\;=\;
    \mathtt{name}\,\mapsto\,
      \llbracket\,
        p_1 \mapsto \varnothing,\ p_2 \mapsto \varnothing,\ \ldots,\
        \varphi \mapsto \llbracket e \rrbracket
      \,\rrbracket.
\end{equation*}
The auto-named variant \ff{> [params] >>} replaces \ff{name} by the
  generated name (\cref{sec:syntax:names});
  the translation is otherwise identical.

\subsection{Meta Directives}\label{sec:sem:meta}

The meta directives recognised by the EO toolchain are listed in
  \cref{tab:metas}.
Meta directives do not contribute \(\varphi\)-expressions to the
  translated program in the same way ordinary objects do;
  they configure name resolution, attach metadata, or instruct downstream
  tooling.

\begin{table*}
\capt{Meta directives recognised at the head of an EO source file.}
\label{tab:metas}
\small
\begin{tabularx}{\textwidth}{l X}
\toprule
Directive & Meaning \\
\midrule
\ff{+package} \(p\)         & Declares the dotted package path of the program. Every top-level object defined in the file is treated as a child of \(\Phi.p\). \\
\ff{+alias} \(p\)           & Imports the canonical object at path \(\Phi.p\) and binds its final segment as an unqualified name in the source (\cref{sec:sem:alias}). \\
\ff{+rt} \(\mathit{platform}\) \(\mathit{coords}\) &
                              Declares a runtime dependency. \(\mathit{platform}\) is an identifier such as \ff{jvm} or \ff{node}; \(\mathit{coords}\) is the platform-specific dependency identifier. \\
\ff{+architect} \(\mathit{email}\) &
                              Author contact for the file. \\
\ff{+home} \(\mathit{url}\) & Repository or project homepage. \\
\ff{+version} \(v\)         & Source version; conventionally a semantic version. \\
\ff{+spdx} \(\mathit{header}\) &
                              An SPDX licence or copyright header preserved verbatim. \\
\ff{+decorate} \(a\) \(b\)  & Declares that the object \(a\) defined in the file is intended to decorate the canonical object at path \(b\). \\
\ff{+unlint} \(\mathit{rule}\) &
                              Suppresses a lint rule for the file. \\
\bottomrule
\end{tabularx}
\end{table*}

Meta directives are legal only at the head of the source, before any
  other object is declared (\cref{sec:syntax:meta});
  the meta header is closed by exactly one blank line.
Unrecognised meta directives are preserved verbatim for downstream
  consumers; the parser does not reject an unknown \ff{+name} at the
  meta-header position.

\subsection{Aliases and Name Resolution}\label{sec:sem:alias}

Identifiers in source come in three forms.
A path rooted at \ff{Q} is absolute:
  \ff{Q.org.eolang.io.stdout} translates to
  \(\Phi.\mathtt{org}.\mathtt{eolang}.\mathtt{io}.\mathtt{stdout}\) with no
  further resolution.
A path rooted at \ff{\$} is scope-relative:
  \ff{\$.x} translates to \(\xi.\mathtt{x}\).
An unqualified identifier (a bare \ff{NAME} as the head of an expression)
  is resolved according to the alias table established by the meta
  header.

The directive \ff{+alias p} imports the canonical object at path
  \(\Phi.p\) and binds the final segment of \(p\) as an unqualified name
  in the current source.
For example, \ff{+alias io.stdout} makes the bare identifier \ff{stdout}
  translate to \(\Phi.\mathtt{io}.\mathtt{stdout}\) wherever it appears
  in the file.
A bare identifier that matches no alias and is not a local attribute
  resolves to \(\xi.\mathtt{name}\) ---
  a dispatch on the current scope --- which is then subject to the
  usual scope-walking rules of the \(\varphi\)-calculus.

The \ff{+package} directive contributes the implicit prefix for the
  file's top-level objects:
  with \ff{+package org.eolang.examples}, a top-level \ff{[] > app}
  contributes the binding
  \(\Phi.\mathtt{org}.\mathtt{eolang}.\mathtt{examples}.\mathtt{app}\,
    \mapsto\,\llbracket\mathtt{[]\ >\ app}\rrbracket\).
